% Response to the qualifying committee
% Pedro C. Campelo Albuquerque, Daniel O. Cajueiro, Rafael T. Menezes
% Submitted: June 2026
\documentclass[11pt]{article}
\usepackage[margin=1in]{geometry}
\usepackage{amsmath,amssymb}
\usepackage{booktabs,longtable}
\usepackage[hidelinks]{hyperref}
\usepackage{xcolor}
\usepackage{enumitem}
\usepackage{microtype}
\usepackage{titlesec}
\titleformat{\section}[hang]{\normalfont\bfseries\large}{\thesection}{0.6em}{}
\setlist[itemize]{leftmargin=1.4em,itemsep=0.3ex,topsep=0.4ex}

\newcommand{\committee}[1]{\textbf{[#1]}}
\newcommand{\page}[1]{(p.~\ref{#1})}

\title{Response to the qualifying committee \\
\large \emph{When Pork-Barrel Backfires: Coalition Presidentialism, Polarization,
and the Price of Legislative Support in Brazil}}

\author{Pedro C. Campelo Albuquerque \and Daniel O. Cajueiro \and Rafael T. Menezes}

\date{June 2026}

\begin{document}

\maketitle

\section*{Overview}

This document summarizes how each of the feedback items raised by the qualifying
committee (Daniel Cajueiro, Bernardo Mueller, and Rafael Terra, defense of
02 June 2026) was addressed in the post-defense revision. We organize the
responses around the four blocks of the internal \texttt{MUSTDO\_v2.md}
roadmap: empirical extensions (Block A), writing improvements (Block B),
external-data integration (Block C), and direct responses to committee
questions (Block D). Each entry indicates the source of the comment
(committee member or our own notes), the original concern, and the
implementation in the revised manuscript. The final draft is
\texttt{docs/paper.pdf} in this repository.

% =====================================================================
\section{Block A — Empirical extensions}
% =====================================================================

\subsection*{A.1 \quad Polarization terciles separately by legislature}
\committee{Bernardo Mueller; Question 3 from the committee}

\noindent\emph{Comment:} ``Terciles in the pooled analysis do not allow us to
distinguish `polarization causes backfire' from `the Bolsonaro era is the phenomenon'.''

\noindent\emph{Response:} The tercile analysis has been re-cut within each
legislature in the revised draft (Subsection~6.3, Table~10). The
within-legislature pattern is informative: the Weak Divergence measure
produces a U-shape that survives the legislature split, and the Strong
and Euclidean measures replicate the cross-legislature gap inside
Legislature 56 alone. The pooled tercile result reported in the
qualifying-defense draft has been superseded by the within-legislature
estimates and is no longer the basis for the central interpretation.

\subsection*{A.2 \quad Structural break in February 2021 (Lira presidency)}
\committee{Our notes, in light of committee discussion}

\noindent\emph{Comment:} ``Under Bolsonaro there is a massive expansion of the
secret budget and Arthur Lira is elected Chamber president. This changes
everything.''

\noindent\emph{Response:} The structural break in February 2021 is documented
explicitly in the revised draft as the \emph{Legislative Capture}
subsection (Subsection~6.4, Table~11, Figure~5). The Lira sub-period
analysis is the empirical core of the post-defense reframing: pork did
not stop functioning under Bolsonaro; the principal of the bargain
shifted from the Executive to the Centrão-led Chamber leadership.
Three sub-samples (Maia~55, Maia~56, Lira~56) and two exclusion
sub-samples (excl.~PP, excl.~Centrão) localize the institutional
transfer.

\subsection*{A.3 \quad Centrão variable and heterogeneity}
\committee{Bernardo Mueller}

\noindent\emph{Comment:} ``Maybe the effect you attribute to polarization is
in fact the Centrão swallowing the amendments.''

\noindent\emph{Response:} The Centrão indicator and the corresponding
heterogeneity analysis are reported in Subsection~6.5 (Table~12). The
sub-sample analysis confirms that the post-Lira Centrão alignment
effect is concentrated within Centrão deputies themselves; once
Centrão deputies are removed from the post-Lira sub-sample, the
estimate is reduced to non-significance. The Centrão composition list
follows the standard nine-party categorization of Power (2010)
and is documented in the table notes.

\subsection*{A.4 \quad Alternative outcome: alignment with the Centrão}
\committee{Our notes}

\noindent\emph{Response:} Implemented in Subsection~6.5 (Table~12). The
exercise was generalized in two directions during the post-defense
revision: alignment with the broader Centrão (the original A.4) and
alignment with the party of the Chamber president (the new
Subsection~6.4, Table~11). The two specifications jointly support the
Legislative Capture interpretation.

\subsection*{A.5 \quad Quadratic treatment specification}
\committee{Committee question 3}

\noindent\emph{Response:} The quadratic specification is implemented as a
robustness exercise in Appendix~F (Table~14). The original Double
Machine Learning estimator with two simultaneously endogenous
variables encountered numerical instability arising from the
near-collinearity of $T$ and $T^2$ in the ElasticNet nuisance
regression; we therefore implemented the augmented model via classical
IV-2SLS with a reduced control set (eleven covariates surviving
QR-based rank validation), within-deputy demeaning, and
deputy-clustered standard errors. The estimates indicate a mild
U-shaped response in all three samples with the inflection located
near \$ R\$3 million per deputy-vote observation. Because the median
$T$ is approximately R\$0.8 million and the 90th percentile is below
R\$3 million, the curvature is observed only in the upper tail and
does not affect the qualitative interpretation of the linear estimates
in Table~4.

\subsection*{A.6 \quad Two-round PEC analysis}
\committee{Bernardo Mueller}

\noindent\emph{Response:} The within-PEC analysis is reported in Appendix~G
(Table~15). We identified nine constitutional amendments with both
first-round and second-round votes recorded in our panel and estimated
the within-PEC change in alignment as a function of the within-PEC
change in committed amendment value, with PEC fixed effects. The
sample is small (3{,}383 deputy-PEC observations), and the precision
of the estimates is limited; we report the exercise as a within-vote
corroboration of the broader pattern rather than as an independent
identification strategy.

\subsection*{A.7 \quad Stylized facts and aggregate descriptives}
\committee{Bernardo Mueller; Our notes}

\noindent\emph{Comment:} ``You have to show the picture: base versus
Centrão, who receives more in proportional terms.''

\noindent\emph{Response:} The new Subsection~2.4 (``Aggregate stylized
facts'') consolidates the descriptive material: annual amendment
commitments by primary-result code 2014--2025 (Table~1, with the
post-ADPF~854 substitution from RP-9 to RP-8 and Pix amendments
visible in the row-by-row comparison), the composition stacked-area
chart (Figure~2), and Centrão versus non-Centrão alignment and
allocation patterns by sub-period (Table~2). The subsection draws on
the federal-budget data described in our response to Block C.

\subsection*{A.8 \quad Consolidated diagnostic test table}
\committee{Bernardo Mueller}

\noindent\emph{Response:} Implemented as Table~3 (\texttt{tab:iv\_validation}).
The table reports first-stage $F$, Anderson--Rubin 95\% CI,
Sargan--Hansen $J$ ($p$-value), and Cinelli--Hazlett $\text{RV}_{q=1}$
in a single panel separated by legislature, allowing integrated
assessment of identification strength.

\subsection*{A.9 \quad Weak vs.\ Strong Divergence as mediator}
\committee{Our notes; companion polarization paper}

\noindent\emph{Response:} The dual-mediation analysis is reported in
Appendix~H (Table~16). The two-mediator Acharya--Blackwell--Sen system
shows that the Strong Divergence channel transmits 31.4\% of the
Temer-era effect and the Weak Divergence channel transmits a
negligible share. The exercise corroborates the single-mediator
analysis in the main text and isolates Strong Divergence as the
relevant transmission mechanism, consistent with the
credibility-cost interpretation developed in Section~7.

\subsection*{A.9.5 \quad Legislative offices as a parallel pork channel}
\committee{Post-defense discussion}

\noindent\emph{Response:} The legislative-office heterogeneity analysis is
reported in Subsection~6.6 (Table~13). Deputies holding Tier-2
positions (party leadership, committee presidency or vice-presidency,
formal rapporteur appointment) display a Bolsonaro-era backfire
sharply attenuated relative to rank-and-file deputies. The pattern is
the observational counterpart of the Legislative Capture
interpretation: deputies with access to procedural levers are insulated
from the credibility cost that drives the rank-and-file backfire.

% =====================================================================
\section{Block B — Writing improvements}
% =====================================================================

\subsection*{B.1 \quad Following the literature's conventions}
\committee{Committee question 1}

\noindent\emph{Comment:} ``On page 10 you have to follow the template of the
literature more closely.''

\noindent\emph{Response:} The empirical-strategy and identification sections
have been rewritten with the conventions of recent Public Choice and
EJPE papers (notation, ordering of identification arguments,
placement of diagnostic tests). The specific page-10 concern is
flagged for a final pass with Bernardo before submission.

\subsection*{B.2 \quad Move descriptive text into tables and figures}
\committee{Committee comment}

\noindent\emph{Response:} The revised draft replaces a substantial portion of
the previously prose-based descriptive content with tables (1--16,
with seven new tables added post-defense) and figures (1--6). The
seasonality figure, the legislative-capture forest plot
(Figure~5), and the quadratic-response curve (Figure~7) are
illustrative examples.

\subsection*{B.3 \quad Justification of the control set}
\committee{Bernardo Mueller}

\noindent\emph{Response:} Section~4 contains a new paragraph documenting the
142-variable full-clean control set, grouped into five families, with
the justification for the wide set resting on (i) the
sparsity-after-selection property of DML, (ii) the multidimensional
Brazilian institutional environment, and (iii) the ablation experiment
reported in Appendix~B (Table~B.13) showing that the structural
coefficient is essentially flat across alternative specifications.

\subsection*{B.4 \quad Reformulating the backfire interpretation}
\committee{Our notes}

\noindent\emph{Response:} Section~7 (Discussion) was substantially rewritten.
The reformulation has three layers: (i) the IV strategy rules out the
reverse-causality interpretation by construction; (ii) the
credibility-cost mechanism (a deputy receiving a visible amendment
incurs a public-legibility cost under polarized constituencies); and
(iii) the Legislative Capture finding shows that the substantive
content of pork-for-votes survives but is re-routed through a
different institutional actor.

% =====================================================================
\section{Block C — External data integration}
% =====================================================================

\subsection*{C.1 \quad Federal-budget aggregates and category composition}
\committee{Committee question 2}

\noindent\emph{Response:} Table~1 of the revised paper reports the
2014--2025 federal-budget aggregates by primary-result category,
sourced from the Federal Transparency Portal and cross-validated
against the Federal Treasury open-data CKAN repository.

\subsection*{C.2 \quad Mandatory execution legal framework}
\committee{Committee question 3}

\noindent\emph{Response:} The new Subsection~2.1 (``The legal architecture
of mandatory execution'') discusses Constitutional Amendments 86 of
2015, 100 of 2019, and 105 of 2019 explicitly.

\subsection*{C.3 \quad Deputy-level RP-9 (Secret Budget) data}
\committee{Committee question 4}

\noindent\emph{Response:} The deputy-level RP-9 indicator $d_{\mathrm{rp9}}$
was constructed from the public list of \emph{apoiamentos}
released under ADPF 854 (Ofícios 932/2024 and 114/2025) and is
incorporated as a control in the augmented specifications reported in
Table~9 (the dual-outcome IV-DML estimates with and without
opaque-channel proxies). An imputed RP-9 measure based on the
prefeito-level GitHub data (Proxy~P4 in
\texttt{STATE\_OF\_PLAY.md}~\S1.5.bis) is used in the legislative-capture
discussion in Subsection~6.4. The official disclosure is
incomplete, and we acknowledge this limitation explicitly in Section~7.

\subsection*{C.4 \quad Post-STF substitution: RP-8 and Pix amendments}
\committee{Our notes}

\noindent\emph{Response:} The post-2022 substitution dynamics from RP-9 to
RP-8 and Pix amendments are documented in Subsections~2.2 and 2.4
(Table~1, Figure~2). The Transparência Brasil et al.~(2025)
technical note on continued ADPF 854 violations is cited in
Sections~2.2 and 7.

% =====================================================================
\section{Block D — Direct responses to committee questions}
% =====================================================================

\subsection*{D.1 \quad Reverse causality}

\noindent\emph{Question:} ``Can you test for reverse causality?''

\noindent\emph{Response:} The IV strategy rules out reverse causality by
construction (the orthogonalization step removes variation in
$T_{it}$ correlated with deputy-specific unobservables). We additionally
report placebo tests (random Bernoulli outcome, post-vote amendment
treatment) in Section~5, all of which produce coefficients indistinguishable
from zero.

\subsection*{D.2 \quad ADPF 850 (sic ADPF 854) as a natural experiment}

\noindent\emph{Question:} ``Do you exploit the natural experiment provided
by the STF intervention?''

\noindent\emph{Response:} The event study around the November 2021 STF
preliminary injunction on RP-9 is reported in Appendix~A
(\texttt{app:event}, Figure~6). We deliberately use the preliminary
injunction rather than the final December 2022 ruling to avoid
confounding the budget intervention with the simultaneous change of
federal administration on 1 January 2023. The event study is
descriptive (no exogenous control group exists) and reported as
supplementary evidence rather than as causal identification.

\subsection*{D.3 \quad Exclusion of Legislature 57}

\noindent\emph{Question:} ``Why do you exclude Legislature 57?''

\noindent\emph{Response:} Two reasons, documented in Subsection~2.2.
First, the official amendment-execution data from the Federal
Transparency Portal exhibits a structural lag of one to two years;
the last fully-closed budget exercise at the time of writing is
2024. Second, Legislature~57 falls inside the post-ADPF~854
substitution regime (RP-8 and Pix amendments), which is
institutionally distinct from the RP-9 era and would mix
heterogeneous regimes in the panel.

\subsection*{D.4 \quad Typology of pork-barrel regimes}

\noindent\emph{Question:} ``Do your findings support a typology of regimes?''

\noindent\emph{Response:} Yes. The Introduction (Section~1) and the
Discussion (Section~7) propose a two-regime typology:
\textbf{Regime~A} (programmatic, cooperative-coalition; Temer
2016--2018) and \textbf{Regime~B} (populist-polarized,
legislatively-captured; Bolsonaro 2019--2022). The two regimes
differ in the sign of the visible-amendment coefficient, in the
identity of the principal of the bargain, and in the salience of
opaque parallel channels.

% =====================================================================
\section*{Summary}
% =====================================================================

Of the 22 substantive items in the post-defense roadmap, 22 have been
implemented in the revised manuscript. The remaining item (B.1, a
final stylistic pass against the literature's page-10 conventions) is
flagged for a last review with Bernardo before submission to Public
Choice. The post-defense reframing reorganizes the central narrative
around the Legislative Capture finding (Subsection~6.4): the
pork-for-votes coefficient does not disappear under populist
polarization; it changes hands. We are grateful for the committee's
careful reading and for the substantive direction the post-defense
feedback provided.

\end{document}
